\documentclass[a4paper,10pt]{article}
\usepackage[utf8]{inputenc}
\usepackage{graphicx}
\usepackage{subfigure}
\usepackage{fancyhdr}
\usepackage{listings}
\usepackage{ gensymb }


%opening

\pagestyle{fancy} %eigener Seitenstil
\fancyhf{} %alle Kopf- und Fußzeilenfelder bereinigen
\fancyhead[L]{SEP Abgabegespräch SS2015/16} %Kopfzeile links
\fancyhead[C]{} %zentrierte Kopfzeile
\fancyhead[R]{\includegraphics[width=1cm]{../../template/tugraz_bw}} %Kopfzeile rechts
\renewcommand{\headrulewidth}{0.4pt} %obere Trennlinie
\fancyfoot[C]{} %Seitennummer
\renewcommand{\footrulewidth}{0.4pt} %untere Trennlinie


\lstset{language=C++, basicstyle=\ttfamily\tiny}  

\begin{document}

\section*{Einfache Datenbank}

Entwickle eine einfache Datenbank in Form eines Key-Value Stores. Das Programm kann Werte mit verschiedenen Datentypen speichern, arbeiten und sie wieder ausgeben. Dafür soll es 3 Befehle geben.

\begin{description}
 \item[set $<$name$>$ $<$type$>$ $<$value$>$] Speichert eine Variable vom Typ \textit{type} in der Datenbank. \textit{type} kann dabei \texttt{``Integer''} oder \texttt{``String''} sein. Ist ein Wert bereits vorhanden, wird er überschrieben.
 \item[get $<$name$>$] Zeigt den Datentyp und Wert der Variable an. Existiert die Variable nicht, soll der Befehl ignoriert werden.
 \item[add $<$name1$>$ $<$name2$>$] Addiert zur Variable \textit{name1} den Wert der Variable \textit{name2}, falls beide vom Typ \texttt{``Integer''} sind. Sind beide vom \texttt{``String''}, sollen die Strings aneinander gehängt werden. In allen anderen Fällen soll ``Incompatible datatypes'' ausgegeben werden. 
\end{description}

Bei der Eingabe von \texttt{bye} oder EOF soll das Programm beendet werden. Der Datentyp \texttt{``Integer''} muss nur Ganzzahlen speichern können, keine Kommazahlen.

\subsection*{Beispiel}
\begin{verbatim}
> set a Integer 3
> set b Integer 2
> set hello String "Hello"
> set world String "World"
> get a
Integer: 3
> set b 4
> add a b
> get a
Integer: 7
> add hello world
> get hello
String: HelloWorld
> get world
String: World
> add hello a
Incompatible datatypes
\end{verbatim}

\subsection*{Aufruf}
\texttt{./db}

\subsection*{Bewertung}
\begin{itemize}
 \item Dokumentation und Programmierstil: 1 Punkt
 \item Einhalten der OOP Konzepte: 1 Punkt
 \item Struktur \& Korrektheit: 1 Punkt
\end{itemize}

\newpage


\end{document}
